\section{Resultados Preliminares, Discussão e Ameaças à Validade}

\subsection{Diagnóstico Empírico das Linhas de Base (Ro\ss kothen et al., 2026)}
Antes da execução dos ciclos com JEV, reprocessamos os 996 registros oficiais ($249 \times 4$ configurações) disponibilizados por Ro\ss kothen et al.~\cite{rosskothen2026zenodo} para estabelecer a linha de base quantitativa do estudo. A Tabela~\ref{tab:baseline_paper} sintetiza o desempenho global e estratificado pelo fator de divisão humano (\textit{Split Factor} -- SF).

\begin{table}[t]
\centering
\caption{Desempenho das Configurações de Referência de Ro\ss kothen et al.~\cite{rosskothen2026zenodo} sobre os 249 Requisitos Fonte}
\label{tab:baseline_paper}
\resizebox{\columnwidth}{!}{%
\begin{tabular}{lcccccc}
\toprule
Configuração & \texttt{CorpBERT} $F_1$ & EMD & MASTeR $F_1$ & Cond. $F_1$ & BLEU & Composto \\
\midrule
\texttt{minimal-qwen235} (Sem Juiz) & \textbf{0,834} & \textbf{0,770} & 0,561 & \textbf{0,758} & \textbf{0,460} & \textbf{0,811} \\
\texttt{full-qwen235} (LLM-Judge)   & 0,823          & 0,762          & 0,553 & 0,742          & 0,423          & 0,797 \\
\texttt{minimal-qwen3-next}         & 0,822          & 0,758          & \textbf{0,592} & 0,725 & 0,412          & 0,801 \\
\texttt{full-gpt-5} (LLM-Judge)     & 0,787          & 0,725          & 0,569 & 0,656          & 0,273          & 0,753 \\
\midrule
\multicolumn{7}{l}{\textit{Degradação em \texttt{full-qwen235} conforme a complexidade estrutural (Split Factor):}} \\
\quad SF $= 1$ ($n=168$)            & 0,864          & 0,785          & 0,601 & ---            & ---            & --- \\
\quad SF $= 2$ ($n=43$)             & 0,766          & 0,734          & 0,478 & ---            & ---            & --- \\
\quad SF $\ge 3$ ($n=38$)           & 0,708          & 0,691          & 0,427 & ---            & ---            & --- \\
\bottomrule
\end{tabular}%
}
\end{table}

Os dados da Tabela~\ref{tab:baseline_paper} confirmam empiricamente as duas hipóteses motivadoras deste estudo:
\begin{enumerate}
    \item \textbf{Evidência de \textit{over-editing} no juiz generativo:} Ao comparar \texttt{minimal-qwen235} com \texttt{full-qwen235}, observa-se que a ativação do loop com juízes LLM generativos reduziu ligeiramente todas as métricas automáticas principais (ex.: escore composto caiu de $0{,}811$ para $0{,}797$), apesar de multiplicar o número de chamadas de inferência. Isso ocorre porque validadores não calibrados emitem críticas estilísticas sobre requisitos simples que já estavam corretos na primeira geração;
    \item \textbf{Colapso em requisitos multi-atômicos:} Em requisitos compostos com $\text{SF} \ge 3$, o \texttt{MasterTemplateF1} despenca de $0{,}601$ para $0{,}427$ ($-28{,}9\%$), demonstrando a necessidade crítica de um mecanismo calibrado de escalonamento humano (\textit{Path C} da nossa proposta).
\end{enumerate}

\subsection{Avaliação Comparativa do JEV-Cascade (RQ1--RQ3)}
A Tabela~\ref{tab:jev_results} consolida os resultados obtidos pelas execuções controladas do repositório experimental (\texttt{OneShot}, \texttt{LLM-Judge} e \texttt{JEV-Cascade}).
% TODO: Após rodar `scripts/run_experiment.py` e `scripts/evaluate_and_export_latex.py`, os números reais serão injetados automaticamente na tabela abaixo e discutidos nos parágrafos de RQ1, RQ2 e RQ3.

\begin{table}[t]
\centering
\caption{Comparação Experimental entre One-Shot, Loop LLM-as-a-Judge e JEV-Cascade ($N = 249$)}
\label{tab:jev_results}
\resizebox{\columnwidth}{!}{%
% Tabela inicial de espaço reservado (sobrescrita automaticamente por scripts/evaluate_and_export_latex.py)
\begin{tabular}{lcccccc}
\toprule
Configuração & MASTeR Main $F_1$ & Cond. $F_1$ & BLEU & Lat. Val. (ms) & Custo/1k (\$) & Escalon. (\%) \\
\midrule
\texttt{gemini\_oneshot}      & \textit{--} & \textit{--} & \textit{--} & $0{,}0$   & \textit{--} & $0{,}0\%$ \\
\texttt{gemini\_llm\_judge}   & \textit{--} & \textit{--} & \textit{--} & \textit{--} & \textit{--} & $0{,}0\%$ \\
\texttt{gemini\_jev\_cascade} & \textit{--} & \textit{--} & \textit{--} & \textit{--} & \textit{--} & \textit{--} \\
\texttt{qwen235\_jev\_cascade}& \textit{--} & \textit{--} & \textit{--} & \textit{--} & \textit{--} & \textit{--} \\
\bottomrule
\end{tabular}

}
\end{table}

\textbf{Análise esperada e guia de interpretação das RQs:}
\begin{itemize}
    \item \textbf{Sobre a RQ1 (Eficácia):} Avalia-se se o \texttt{JEV-Cascade} preserva os altos índices de \texttt{OneShot} nos itens aprovados via \textit{Fast Accept} enquanto recupera falhas gramaticais estritas (ausência de verbo modal ou prefixo condicional incorreto) por meio do mapeamento determinístico de reparo;
    \item \textbf{Sobre a RQ2 (Economia e Latência):} Diferentemente do modo \texttt{LLM-Judge}, no qual cada candidato aciona geração autorregressiva de comentários e uma etapa adicional de filtragem, cada auditoria JEV executa em um único passe paralelo sem tokens de saída, reduzindo substancialmente a latência acumulada de validação e o custo por milhar de requisitos;
    \item \textbf{Sobre a RQ3 (Calibração da Cascata):} Compara-se o \texttt{MasterTemplateF1} médio do subconjunto aceito autonomamente contra o subconjunto escalado ($\Delta F_1$), verificando se a zona de incerteza ($\tau_{\text{rej}} \le p_j < \tau_{\text{acc}}$) concentra efetivamente os casos com $\text{SF} \ge 2$.
\end{itemize}

\subsection{Ameaças à Validade}
Seguindo as diretrizes metodológicas de Runeson e H\"ost~\cite{runeson2009guidelines}, destacamos quatro categorias de ameaças à validade e suas respectivas mitigações:
\begin{itemize}
    \item \textbf{Validade de Construto:} As métricas automáticas comparam as formalizações geradas contra uma única reescrita humana de referência~\cite{grosser2023comparative}. Como demonstrado por Ro\ss kothen et al.~\cite{rosskothen2026converting}, existem múltiplas maneiras válidas de decompor requisitos complexos no padrão MASTeR. Mitigamos essa ameaça combinando métricas estruturais de multiconjunto (\texttt{MasterTemplateF1}) com métricas léxico-semânticas e analisando separadamente os estratos de complexidade;
    \item \textbf{Validade Interna:} Alterar simultaneamente o LLM gerador e o mecanismo validador impediria isolar a contribuição do JEV. Mitigamos esse risco preservando integralmente os prompts Jinja2 originais~\cite{rosskothen2026zenodo} e executando os três modos (\texttt{OneShot}, \texttt{LLM-Judge} e \texttt{JEV-Cascade}) dentro da mesma família de gerador (\texttt{Gemini 2.5 Flash}), além de suportar execução direta com \texttt{Qwen3-235B};
    \item \textbf{Validade Externa:} O estudo concentra-se no sistema de templates MASTeR sobre 249 requisitos de cinco projetos (predominantemente sistemas críticos e embarcados). Embora representativo de cenários industriais como PREEvision, a generalização para outros padrões semiformais (como EARS) ou requisitos ágeis (\textit{User Stories}) requer investigação adicional;
    \item \textbf{Confiabilidade:} Modelos em nuvem podem sofrer variações de latência de rede e atualizações de versão. Para garantir reprodutibilidade estrita, fixamos temperatura igual a zero ($T = 0{,}0$) com decodificação estruturada, persistimos telemetria individual por chamada em arquivos JSONL e disponibilizamos todo o pipeline de avaliação e testes automatizados.
\end{itemize}
